% Generated draft: requires mathematical and source review.
\documentclass[11pt]{article}
\usepackage[margin=25mm]{geometry}
\usepackage{fontspec}
\setmainfont{Latin Modern Roman}
\usepackage{amsmath,amssymb,mathtools,unicode-math}
\setmathfont{Latin Modern Math}
\usepackage{enumitem,needspace,xurl,hyperref}
\setlength{\parindent}{0pt}
\setlength{\parskip}{5pt}
\setlength{\emergencystretch}{3em}
\newcommand{\N}{\mathbb{N}}
\newcommand{\Z}{\mathbb{Z}}
\newcommand{\Q}{\mathbb{Q}}
\newcommand{\R}{\mathbb{R}}
\newcommand{\C}{\mathbb{C}}
\newcommand{\F}{\mathbb{F}}
\newcommand{\A}{\mathbb{A}}
\newcommand{\PP}{\mathbb{P}}
\newcommand{\E}{\mathbb{E}}
\newcommand{\eps}{\varepsilon}
\newcommand{\dd}{\,\mathrm{d}}
\begin{document}
\section*{UnsolvedMath ID: 30006478}
\textbf{Source problem code:} catalogue entry
\subsection{Definitions and mathematical statement}
The original source attributes the question to Michael Stoll and uses the term ``nice curves.'' Under the usual interpretation of this term, a curve here is a smooth, projective, geometrically integral curve over \(\Q\). The cleaned registry statement explicitly specifies smoothness and projectivity but does not explicitly specify geometric integrality. The formulation below uses the usual interpretation, subject to confirmation of this unresolved convention.

For such a curve \(C\), let \(C(\Q)\) denote its set of \(\Q\)-rational points, let \(\#C(\Q)\) denote its cardinality, and let
\[
 g(C)=\dim_{\Q}H^1(C,\mathcal O_C)
\]
denote its genus.

Determine whether there exists a sequence \((C_n)_{n\geq1}\) of such curves over \(\Q\) satisfying
\[
 g(C_n)\geq2\quad\text{for every }n\geq1,
 \qquad
 \lim_{n\to\infty}\frac{\#C_n(\Q)}{g(C_n)}=+\infty.
\]
Explicitly, the limit condition means
\[
 \forall M\in\R_{>0}\ \exists N\in\N_{\geq1}\ \forall n\geq N,
 \qquad \#C_n(\Q)>M\,g(C_n).
\]
The requested task is to construct such a sequence or prove that no such sequence exists. No condition \(g(C_n)\to\infty\) is imposed in the question itself.

The supplied original remark states that the weak Lang conjecture implies a bound on the number of rational points depending only on the genus, and hence would force any sequence satisfying the displayed limit to have \(g(C_n)\to\infty\). It also states that a bound in terms of genus and Mordell--Weil rank has been established; consequently, under a uniform bound on the Mordell--Weil ranks of the Jacobians of the curves in the sequence, the same conclusion \(g(C_n)\to\infty\) would follow. These are supplied contextual claims, not additional hypotheses of the question or independently verified assertions here. The original references labeled [3] and [4] are not bibliographically identified in the supplied excerpt.
\subsection{Short English statement}
Can one find smooth projective curves over \(\Q\), each of genus at least two, whose numbers of rational points divided by their genera tend to infinity, or prove this impossible? The original term ``nice'' is interpreted here as also requiring geometric integrality, a convention requiring confirmation. Genus tending to infinity is only a consequence discussed under additional assumptions in the supplied remarks.
\subsection{Sources}
\begin{itemize}
\item Supplied UnsolvedMath local registry record, permanent ID \(30006478\), title ``Curves with Unbounded Rational-Point-to-Genus Ratio,'' exact source problem code \texttt{catalogue entry}. The original statement attributes the question to Michael Stoll; the registry records proposed year 2025. The registry labels the problem open and the statement \texttt{corrected\_verified}, explaining that only formula delimiters, notation, grammar, and typography were repaired. These are supplied registry labels, not independent verification.
\item Supplied report citation: \emph{Arithmetic Statistics for Algebraic Objects (2026)}, Oberwolfach Reports, DOI 10.4171/owr/2025/52, pp. 2765--2842. \url{https://doi.org/10.4171/owr/2025/52}. The literature-source list calls this ``Oberwolfach report 52/2025.'' The differing year descriptions are retained without resolution.
\item Supplied literature reference: Dimitrov--Gao--Habegger, ``uniformity in Mordell--Lang.'' \url{https://annals.math.princeton.edu/2021/194-1/p05}. The supplied excerpt does not explicitly identify this link with its reference [4].
\item Literature assessment checked at \texttt{2026-08-22}: the supplied curation labels the problem open and states that uniformity results bound rational points in important low-rank or conditional settings, but that no construction or impossibility theorem is known for the requested sequence over \(\Q\). This is supplied dated curation, not independently verified. Its instructions call for independent verification of time-sensitive status and cited claims.
\end{itemize}
\end{document}
